\scanpage{41}
hom de \deleted{\unsure{fais}} foncteurs en groupes sur $(\mathrm{Sch})_{/S}$] on peut associer à $\xi$ un foncteur $F_\xi$, en utilisant $\unsure{\delta_n|\mathscr D_n}$.

On veut caractériser les foncteurs obtenus de cette façon. Notons que $\xi$ se \deleted{\unsure{trouve. vice}} \inserted{retrouve} à partir de $F_\xi$ en \deleted{\unsure{lui faisant}} \deleted{\inserted{faisant \unsure{opé}}} \inserted{\unsure{intégrant|interprétant}} \deleted{$S' = S,\ X\quad X'=E_S^n$ sur $S'$,} [\unsure{au dessus|en dessous}

\diagram{\begin{tikzcd}[ampersand replacement=\&,column sep=small,row sep=tiny,/tikz/execute at end picture={\draw[thin] (current bounding box.north west) -- (current bounding box.south east);}]
 \deleted{\xi_{E^n}} \& \deleted{\xi_{E^n_{S^{d,n}}}} \& \deleted{\xi} \& \\
 \& \deleted{\unsure{S'}\quad S_n=} \& \& \deleted{X=E_S^n}\\
 \deleted{E^n}\arrow[d] \& \deleted{E^n_{S^{d,n}}}\arrow[l]\arrow[d] \& \deleted{E^n}\arrow[l]\arrow[d] \& \deleted{S}\arrow[l]\\
 \deleted{S} \& \deleted{S^{d,n}}\arrow[l]\arrow[ul,dashed] \& \deleted{S}\arrow[l] \&\\
 \& \deleted{\Vert}\&\&\\
 \& \deleted{\operatorname{Spec}\mathcal O_S[t_1,\ldots t_d]/(t_1,\ldots t_d)^{n+1}}\&\&
\end{tikzcd}}

\deleted{de $S$ pour simplifier, mais \unsure{on répète|en le faisant} sur tt \unsure{$S'$|$T'$}]} \deleted{$G^{d,\unsure{n}}_{\deleted{\unsure{S}}}(S)$ comme $\underline{\operatorname{Aut}}_S(S^{d,n})$, que}
\inserted{\deleted{[augm $\varepsilon:S\longrightarrow S^{d,n}$]}}
\deleted{on \unsure{donne|laisse} opérer sur $E^n_{S^{d,n}}$ en laissant fixe $E^n$ considéré comme $E^n_{S^{d,n}}\times_{S^{d,n}}S$. Un tel automorphisme opère donc sur $\xi_{E^n_{S^{d,n}}}=\xi_{E^n}\times_S S^{d,n}$, de façon compatible avec les opérations sur $E^n_{S^{d,n}}$,}

% archive sheet number: 40

